% ============================================================================
% ANNEXE 1 -- CONTENEURISATION DE LA PLATEFORME GHABETNA
% ----------------------------------------------------------------------------
% Redigee a partir du code reel du projet :
%   - Docker-compose.yml (racine du depot)
%   - backend/{auth_ms,forest_ms,alert_ms,analytics_ms,chatbot_ms,citizen_ms}/Dockerfile
%   - gateway/Dockerfile
%   - frontend/forest_app/Dockerfile + nginx.conf
%
% Packages requis : aucun au-dela de la classe de base (tabular, verbatim,
% minipage et fbox sont natifs). Si votre preambule charge « float », vous
% pouvez remplacer les [h] par des [H] pour figer la position des tableaux.
% ============================================================================

\newpage
\phantomsection
\addcontentsline{toc}{chapter}{Annexe 1 : Conteneurisation de la plateforme}
\markboth{Annexe 1 : Conteneurisation de la plateforme}{}

\begin{center}
  {\fontsize{30}{35}\selectfont \textbf{Annexe 1 : Conteneurisation de la plateforme}}
\end{center}
\vspace{1cm}

Cette annexe décrit la démarche de conteneurisation de la plateforme Ghabetna :
les motivations qui l'ont justifiée, le périmètre couvert, la structure des
images produites et leur orchestration au sein d'un environnement unique.

\section*{Pourquoi conteneuriser Ghabetna ?}

L'architecture de Ghabetna ne se limite pas à une application unique. Elle
réunit six microservices FastAPI indépendants (\texttt{auth\_ms},
\texttt{forest\_ms}, \texttt{alert\_ms}, \texttt{analytics\_ms},
\texttt{citizen\_ms} et \texttt{chatbot\_ms}), une API Gateway, cinq bases de
données PostgreSQL dont deux étendues par PostGIS, une instance Redis servant
de broker d'événements, ainsi que l'application web Flutter destinée aux
administrateurs et aux superviseurs.

Sans conteneurisation, installer cet ensemble sur une nouvelle machine
supposerait de reproduire manuellement chacun de ces composants : cinq
instances PostgreSQL sur des ports distincts, l'extension PostGIS et ses
bibliothèques système (\texttt{libgdal}, \texttt{libgeos}), un serveur Redis,
puis les dépendances Python propres à chaque service. Le service
\texttt{chatbot\_ms} illustre particulièrement cette difficulté : il repose
sur \texttt{torch} et \texttt{sentence-transformers}, dont l'installation
manuelle est longue et sensible à la version de Python utilisée.

Docker répond à ce besoin en encapsulant chaque service avec l'ensemble de
ses dépendances dans une unité isolée et reproductible. Pour Ghabetna, cela
se traduit concrètement par quatre garanties :

\begin{itemize}
  \item \textbf{Reproductibilité :} la version de Python, celle de PostgreSQL
        et celle de PostGIS sont fixées dans les images. Le comportement
        obtenu est identique entre le poste de développement et tout autre
        environnement, ce qui écarte les écarts de configuration du type
        \textit{« ça fonctionne sur ma machine »}.

  \item \textbf{Isolation des services :} chaque microservice dispose de sa
        propre base de données isolée, conformément au principe
        \textit{database-per-service} retenu au chapitre~2. Les schémas de
        \texttt{auth\_db}, \texttt{forest\_db}, \texttt{alert\_db},
        \texttt{analytics\_db} et \texttt{citizen\_db} ne peuvent ainsi pas
        interférer entre eux.

  \item \textbf{Cohérence des traitements géospatiaux :} les services
        \texttt{forest\_ms} et \texttt{alert\_ms} exécutent des requêtes
        PostGIS (\texttt{ST\_Overlaps}, \texttt{ST\_Contains},
        \texttt{ST\_Area}). L'image \texttt{postgis/postgis:15-3.3} garantit
        que la version de l'extension reste identique partout, condition
        nécessaire pour que les contrôles de chevauchement produisent les
        mêmes résultats d'un environnement à l'autre.

  \item \textbf{Démarrage unifié :} une seule commande
        \texttt{docker compose up -{}-build} lance l'intégralité de la
        plateforme, bases de données, Redis, microservices, API Gateway et
        frontend inclus, dans le bon ordre de dépendance.
\end{itemize}

\section*{Périmètre de la conteneurisation}

L'ensemble du backend ainsi que l'application web Flutter sont conteneurisés,
soit quatorze conteneurs orchestrés par un unique fichier
\texttt{Docker-compose.yml} placé à la racine du dépôt. Le
tableau~\ref{tab:docker-images} récapitule ces conteneurs, leur rôle, leur
image de base et le port exposé sur la machine hôte.

\begin{table}[h]
\centering
\caption{Conteneurs de la plateforme Ghabetna}
\label{tab:docker-images}
\renewcommand{\arraystretch}{1.25}
\begin{tabular}{|p{2.7cm}|p{4.3cm}|p{3.9cm}|p{1.5cm}|}
\hline
\textbf{Conteneur} & \textbf{Rôle} & \textbf{Image de base} & \textbf{Port hôte} \\
\hline
\multicolumn{4}{|l|}{\textbf{Application cliente}} \\
\hline
\texttt{frontend} & Application web Flutter (superviseurs, administrateurs)
& \texttt{flutter:stable} \newline puis \texttt{nginx:alpine} & 8080 \\
\hline
\multicolumn{4}{|l|}{\textbf{Point d'entrée et microservices}} \\
\hline
\texttt{gateway} & Point d'entrée unique, routage et vérification des JWT
& \texttt{python:3.11-slim} & 8000 \\
\hline
\texttt{auth\_ms} & Authentification, rôles et activation de compte
& \texttt{python:3.11-slim} & 8001 \\
\hline
\texttt{forest\_ms} & Forêts, parcelles et affectations (PostGIS)
& \texttt{python:3.11-slim} & 8002 \\
\hline
\texttt{alert\_ms} & Signalement et suivi des incidents géolocalisés
& \texttt{python:3.11-slim} & 8003 \\
\hline
\texttt{analytics\_ms} & Agrégation statistique et KPIs
& \texttt{python:3.11-slim} & 8004 \\
\hline
\texttt{chatbot\_ms} & Assistant réglementaire RAG (bge-m3, BM25, Groq)
& \texttt{python:3.12-slim} & 8005 \\
\hline
\texttt{citizen\_ms} & Inscription des usagers et validation des dossiers
& \texttt{python:3.11-slim} & 8006 \\
\hline
\multicolumn{4}{|l|}{\textbf{Persistance et communication}} \\
\hline
\texttt{auth\_db} & Utilisateurs, rôles et tokens
& \texttt{postgres:15} & 5433 \\
\hline
\texttt{forest\_db} & Polygones des forêts et des parcelles
& \texttt{postgis/postgis:} \texttt{15-3.3} & 5434 \\
\hline
\texttt{alert\_db} & Incidents géolocalisés
& \texttt{postgis/postgis:} \texttt{15-3.3} & 5435 \\
\hline
\texttt{analytics\_db} & Statistiques agrégées
& \texttt{postgres:15} & 5436 \\
\hline
\texttt{citizen\_db} & Dossiers d'inscription et justificatifs
& \texttt{postgres:15} & 5437 \\
\hline
\texttt{redis} & Broker d'événements (Redis Streams)
& \texttt{redis:7-alpine} & 6380 \\
\hline
\end{tabular}
\end{table}

Deux composants du projet restent en dehors de ce périmètre, pour une raison
qui tient à leur mode de distribution plutôt qu'à une limite technique. Les
applications \texttt{agent\_app} et \texttt{public\_user\_app} sont des
applications mobiles Flutter : elles sont compilées en paquets Android
(\texttt{APK}) installés directement sur l'appareil de l'utilisateur, et non
servies par un serveur. Leur conteneurisation n'aurait donc pas d'objet.
Seule l'application \texttt{forest\_app}, consultée depuis un navigateur, est
compilée en fichiers statiques puis servie par un conteneur web.

Le service \texttt{chatbot\_ms} constitue par ailleurs le seul microservice
dépourvu de base de données relationnelle : son index (embeddings et BM25)
est constitué de fichiers locaux \texttt{.npy} et \texttt{.pkl}, montés en
volume plutôt que stockés en base.

\section*{Structure des Dockerfiles}

\subsection*{Microservices FastAPI}

Les sept Dockerfiles Python (API Gateway et six microservices) suivent une
structure commune volontairement minimale : image de base
\texttt{python:3.11-slim}, définition du répertoire de travail, copie du
fichier \texttt{requirements.txt} et installation des dépendances
\textbf{avant} la copie du code source, puis copie du code.

Cet ordre n'est pas arbitraire. Docker met en cache chaque instruction sous
forme de couche : en installant les dépendances avant de copier le code, la
couche d'installation, de loin la plus longue, n'est reconstruite que si le
fichier \texttt{requirements.txt} change. Une simple modification du code
source laisse cette couche intacte, ce qui réduit sensiblement le temps de
reconstruction pendant le développement.

\begin{verbatim}
FROM python:3.11-slim
WORKDIR /app
COPY requirements.txt .
RUN pip install --no-cache-dir -r requirements.txt
COPY . .
CMD ["sh", "-c", "alembic upgrade head && uvicorn app.main:app \
     --host 0.0.0.0 --port 8001"]
\end{verbatim}

La commande de démarrage est définie dans chaque Dockerfile par l'instruction
\texttt{CMD}, et non dans le fichier \texttt{docker-compose.yml}. Elle
enchaîne systématiquement deux opérations : l'application des migrations
Alembic (\texttt{alembic upgrade head}), puis le lancement du serveur uvicorn
sur le port propre au service. Ce couplage garantit qu'un conteneur qui
démarre travaille toujours sur un schéma de base à jour, sans intervention
manuelle après un déploiement. L'option \texttt{-{}-reload} n'est activée
dans aucun service : elle relancerait le processus à chaque modification de
fichier, comportement utile en développement local mais indésirable dans un
conteneur.

\subsection*{Cas particuliers}

Trois services s'écartent de ce modèle commun, chaque écart répondant à une
contrainte identifiée.

\paragraph{Services géospatiaux.}
\texttt{forest\_ms} et \texttt{alert\_ms} manipulent des géométries PostGIS
au moyen de GeoAlchemy2. Leur Dockerfile installe au préalable les
bibliothèques système correspondantes (\texttt{libgdal-dev},
\texttt{gdal-bin}, \texttt{libgeos-dev}) via \texttt{apt-get}, suivi d'un
\texttt{rm -rf /var/lib/apt/lists/*} qui supprime le cache des paquets et
évite d'alourdir inutilement l'image finale.

\paragraph{Attente de la base dans \texttt{alert\_ms}.}
\texttt{alert\_ms} installe en outre \texttt{postgresql-client} afin de
disposer de l'utilitaire \texttt{pg\_isready}. Sa commande de démarrage
attend explicitement que la base réponde avant de lancer les migrations :

\begin{verbatim}
CMD ["sh", "-c", "until pg_isready -h alert_db -p 5432 -U alert_user; \
     do echo 'Waiting for DB...'; sleep 2; done && alembic upgrade head \
     && uvicorn app.main:app --host 0.0.0.0 --port 8003"]
\end{verbatim}

Cette boucle complète le \texttt{healthcheck} déclaré côté Compose : elle
protège le service du cas où la base accepterait les connexions réseau sans
être encore prête à traiter les requêtes, situation qui ferait échouer la
migration au démarrage.

\paragraph{Service conversationnel.}
\texttt{chatbot\_ms} se distingue sur trois points. Il repose sur
\texttt{python:3.12-slim} et installe \texttt{build-essential} ainsi que
\texttt{curl}, nécessaires à certaines dépendances de \texttt{torch}. Il ne
copie ensuite que les répertoires utiles à l'exécution (\texttt{rag/},
\texttt{app/}, \texttt{scripts/}) plutôt que l'intégralité du dossier, ce qui
écarte de l'image le corpus, les scripts d'évaluation et la documentation.
Enfin, le modèle d'embedding \texttt{BAAI/bge-m3}, qui pèse environ 2,2~Go,
n'est pas intégré à l'image : la variable \texttt{HF\_HOME} pointe vers un
répertoire de cache monté en volume nommé (\texttt{hf\_cache}). Sans ce
choix, le modèle serait retéléchargé à chaque reconstruction de l'image.

\subsection*{Frontend Flutter (build multi-étapes)}

Le Dockerfile de l'application web exploite un \textbf{build multi-étapes},
technique qui sépare l'étape de compilation de l'étape de service afin de
produire une image finale ne contenant aucun outil de développement :

\begin{enumerate}
  \item \textbf{Étape 1 (compilation) :} l'image
        \texttt{ghcr.io/cirruslabs/flutter:stable} récupère les dépendances
        Dart (\texttt{flutter pub get}) puis compile l'application en
        fichiers statiques avec \texttt{flutter build web -{}-release}. Les
        fichiers \texttt{pubspec.yaml} et \texttt{pubspec.lock} sont copiés
        avant le reste du code, pour la même raison de cache que le
        \texttt{requirements.txt} des services Python. Cette étape n'est pas
        conservée dans l'image finale : elle sert uniquement d'environnement
        de compilation.

  \item \textbf{Étape 2 (service) :} l'image finale repose sur
        \texttt{nginx:alpine}. Elle copie uniquement le répertoire
        \texttt{build/web/} produit par l'étape précédente, ainsi qu'un
        fichier \texttt{nginx.conf} personnalisé. Le SDK Flutter, la chaîne
        de compilation Dart et les sources ne figurent pas dans l'image
        livrée, dont la taille s'en trouve fortement réduite.
\end{enumerate}

Le fichier \texttt{nginx.conf} configure un serveur adapté à une application
monopage. Sa directive essentielle est
\texttt{try\_files \$uri \$uri/ /index.html} : toute route inconnue du
serveur est redirigée vers \texttt{index.html}, ce qui laisse le routeur
côté client (\texttt{go\_router}) prendre en charge la navigation. Sans
cette règle, l'accès direct à une adresse interne de l'application, ou un
simple rechargement de page, renverrait une erreur 404.

\section*{Orchestration avec Docker Compose}

Le fichier \texttt{Docker-compose.yml} placé à la racine du dépôt décrit
l'ensemble des conteneurs et leurs relations. Ses points de configuration
notables sont les suivants.

\begin{itemize}
  \item \textbf{Variables d'environnement.} Chaque service dispose de son
        propre fichier \texttt{.env}, déclaré par la directive
        \texttt{env\_file}, qui regroupe ses paramètres sensibles :
        \texttt{JWT\_SECRET\_KEY} et identifiants SMTP pour
        \texttt{auth\_ms}, \texttt{GROQ\_API\_KEY} pour \texttt{chatbot\_ms},
        \texttt{SECRET\_KEY} pour la Gateway. Ces fichiers sont exclus du
        dépôt Git par le \texttt{.gitignore} : aucune donnée sensible n'est
        ainsi versionnée ni codée en dur dans les images. Les paramètres non
        sensibles sont en revanche déclarés directement dans le bloc
        \texttt{environment} du fichier Compose, notamment les chaînes de
        connexion aux bases et les URL inter-services.

  \item \textbf{Volumes nommés.} Chaque base de données dispose d'un volume
        dédié (\texttt{auth\_db\_data}, \texttt{forest\_db\_data},
        \texttt{alert\_db\_data}, \texttt{analytics\_db\_data},
        \texttt{citizen\_db\_data}) qui assure la persistance des données
        entre deux redémarrages. Deux volumes supplémentaires,
        \texttt{alert\_uploads} et \texttt{citizen\_uploads}, conservent
        respectivement les photographies jointes aux signalements et les
        pièces justificatives déposées par les usagers : sans eux, ces
        fichiers disparaîtraient à chaque reconstruction du conteneur
        correspondant. Le volume \texttt{hf\_cache} joue le même rôle pour le
        modèle d'embedding du chatbot, tandis que le corpus réglementaire est
        monté depuis le système de fichiers hôte, ce qui permet d'ajouter ou
        de mettre à jour un texte sans reconstruire l'image.

  \item \textbf{Healthchecks et ordre de démarrage.} Les cinq bases de
        données exposent un \texttt{healthcheck} fondé sur
        \texttt{pg\_isready}, et Redis un \texttt{redis-cli ping}. Les
        microservices déclarent leurs dépendances au moyen de
        \texttt{depends\_on}, en distinguant deux conditions. La condition
        \texttt{service\_healthy} est retenue lorsqu'un service ne peut pas
        démarrer sans son composant : \texttt{auth\_ms} attend que
        \texttt{auth\_db} réponde effectivement, puisqu'il applique ses
        migrations dès le lancement. La condition \texttt{service\_started},
        moins stricte, est utilisée lorsqu'un service n'en appelle un autre
        qu'au moment d'une requête utilisateur ; \texttt{citizen\_ms}
        n'interroge par exemple \texttt{auth\_ms} qu'à la première
        inscription, si bien qu'un léger décalage de démarrage reste sans
        effet.

  \item \textbf{Réseau applicatif dédié.} Tous les conteneurs sont rattachés
        à un réseau \texttt{bridge} nommé \texttt{ghabetna\_net}, déclaré
        explicitement plutôt que laissé au réseau par défaut. Ce réseau leur
        permet de se joindre par nom de service
        (\texttt{http://auth\_ms:8001}, \texttt{redis://redis:6379}) en
        s'appuyant sur la résolution DNS interne de Docker. Les ports
        internes restent inchangés quel que soit le port publié sur l'hôte :
        \texttt{redis} est ainsi joignable sur le port \texttt{6379} depuis
        le réseau applicatif, tout en étant exposé sur \texttt{6380} à
        l'extérieur afin de ne pas entrer en conflit avec une éventuelle
        instance Redis déjà installée sur la machine.
\end{itemize}

\section*{Démarrage de la plateforme}

L'ensemble des conteneurs est construit et lancé par une commande unique,
exécutée à la racine du projet :

\begin{verbatim}
docker compose up --build -d
\end{verbatim}

Docker construit alors les images à partir des Dockerfiles, télécharge les
images de base manquantes, crée les volumes et le réseau
\texttt{ghabetna\_net}, puis démarre les quatorze conteneurs en respectant
l'ordre imposé par les \texttt{healthchecks}. L'état des services et leurs
journaux sont ensuite consultables par les commandes suivantes :

\begin{verbatim}
docker compose ps
docker compose logs -f alert_ms
\end{verbatim}

La plateforme est dès lors accessible : l'interface web d'administration sur
\texttt{http://localhost:8080}, l'API Gateway sur
\texttt{http://localhost:8000}, et la documentation OpenAPI auto-générée de
chaque microservice sur son port respectif. L'arrêt s'effectue par
\texttt{docker compose down}, les volumes nommés préservant les données
au-delà de la durée de vie des conteneurs.
